\section*{Levy's Continuity Theorem and Central Limit Theorem}

\subsection*{14.1 Weak Convergence}

14Levy's Continuity Theorem and Central Limit

Theorem

The characteristic function of a random variable is a powerful tool that encapsulates

all the information about its probability distribution. Levy's continuity theorem

establishes the equivalence between pointwise convergence of characteristic func-

tions and convergence in distribution. This theorem offers a convenient way to

determine whether a sequence of random variables converges in distribution, and

serves as a tool for proving the central limit theorem.

\begin{thmbox}{14.1 (Levy's Continuity Theorem)}
\end{thmbox}

Let .(Xn)\infty

n=1 be a sequence of real-valued random variables. Let Pn denote

the probability measure induced by Xn on R, and \phin(t) be the characteristic

function of Xn,f o r n\geq 1 Suppose the characteristic functions converge

pointwise to some function \phi(t).

Then the followings are equivalent:

1. Xn converges weakly to some random variable X as n tends to infinity.

2. \phi(t) is the characteristic function of a random variable.

3. \phi(t) is continuous at t= 0 .

4. The sequence of probability measures .(Pn)\infty

n=1 is tight.

In the first section in this chapter, we define the notion of weak convergence,

which is equivalent to convergence in distribution. We also prove that convergence

in total variation implies convergence in distribution. Then, we introduce the notion

of tightness of measures and Prokhorov theorem. In the last section, we use the

continuity theorem to prove a version of central limit theorem that assumes iid.

random variables with finite variance. In the last section, we will state a central

limit theorem for triangular array and illustrate it through an example.

Suppose we can only obtain information about a probability distribution P through

some measurements of the form .

\int

hdP , where h(x) is a "test function"If we

have a sequence of probability distributions Pn's, it would be desirable if, for each

test function h(x), the sequence of measurements .

\int

hdP n converges as n tends to

infinity. This motivates the concept of weak convergence, which allows us to study

the convergence of probability distributions based on their behavior with respect to

test functions.

\begin{defbox}{14.1}
\end{defbox}

Let. Pn,f o rn \geq 1, be probability measures and P be another probability measure

defined on sample space RI f

limn\to\infty

\int

R

h(x) dPn(x) =

\int

R

h(x) dP(x)

for all bounded and continuous function h(x), then we say that Pn converges

weakly to P The notation for weak convergence is

Pn

W

\to P or Pn P.

We often apply this definition with the probability measures are the push-forward

measures of some random variables. By using the change-of-variable formula

(Theorem 7.11), we can formulate weak convergence for a sequence of random

variables as follows.

\begin{defbox}{14.2}
\end{defbox}

Suppose X1, X2, X3 ... , is a sequence of random variables, and X is a random

variable. We say that .(Xn)\infty

n=1 converges weakly to X if the image measures of

Xn's converge weakly to the image measure of X, ie.,

limn\to\infty E[h(Xn)]= E[h(X)]

for all continuous and bounded functions h.

The next theorem establishes the equivalence of weak convergence and conver-

gence in distribution. Weak convergence is an elegant definition, and it is more

convenient for proofs. However, checking the condition for weak convergence can

be challenging, as it involves considering infinitely many potential test functions. In

contrast, the definition of convergence in distribution is often more practical to use

in calculations.

\begin{thmbox}{14.2 (Weak Convergence and Convergence in Distribution)}
\end{thmbox}

Suppose .(Xn)\infty

n=1 is a sequence of random variables, and X be another

random variable. Denote the cumulative distribution function of. Xn byFn(x),

forn \geq 1, and the cumulative distribution function of X byF(x) We have

Xn

W

\to X if and only if Fn

D

\to F.

\textit{Proof} (. ) Suppose F(x) is continuous at x = a, ie., P({a}) = 0. The idea of

proof is to approximate an indicator function by piece-wise linear function. Define

a piece-wise linear function g by

g(x) =

1i f x<a - \epsilon,

\epsilon(a - x) if a - \epsilon\leq x \leq a,

0i f x>a .

The function g is linearly decreasing from 1 to 0 in the interval .[a - \epsilon,a] and is

continuous and bounded. We have

1(-\infty ,a-\epsilon](x) \leq g(x) \leq 1(-\infty ,a](x).

Substitute x by X and Xn in the above inequalities, and take expectation to get

F(a - \epsilon)\leq E[g(X)]= limn\to\infty E[g(Xn)]\leq lim infn\to\infty Fn(a). (14.1)

The last inequality in (14.1) follows from E[g(Xn)]\leq Fn(a) for all n.

Similarly, consider the function h(x) = g(x - \epsilon)The function h(x) satisfies

1(-\infty ,a](x) \leq h(x) \leq 1(-\infty ,a+\epsilon](x)

and is continuous and bounded. We thus obtain

lim sup

n\to\infty

Fn(a) \leq limn\to\infty E[h(Xn)]= E[h(X)]\leq F(a + \epsilon). (14.2)

Putting (14.1) and (14.2) together, we get

F(a - \epsilon)\leq lim infn\to\infty Fn(a) \leq lim sup

n\to\infty

Fn(a) \leq F(a + \epsilon).

Since F(x) is continuous at x = a, the limit of .(Fn(a))n\geq1 exists and equals F(a) .

(. ) We apply Skorokhod representation theorem (Theorem 10.8), and let

(Yn)n\geq1 be a sequence of random variables defined on the same probability space,

such that Yn and Xn have the same distribution and Yn's converge almost surely to

a random variable Y that has the same distribution as X.

Leth(y) be a continuous and bounded function on. R. Because h is continuous, by

the continuous mapping theorem (Theorem 10.11),h(Yn) converges toh(Y) almost

surely. Since h(Yn) is bounded, we can apply the dominated convergence theorem

to obtain

limn\to\infty E[h(Xn)]= limn\to\infty E[h(Yn)]= E[limn\to\infty h(Yn)]= E[h(Y)]= E[h(X)].

Because it holds for any continuous and bounded function h(x), this proves that Xn

converges weakly to X. \hfill $\square$

We can now prove that Condition 1 implies Condition 2 in Theorem 14.1.

\begin{thmbox}{14.3}
\end{thmbox}

IfXn

W

-\to X , then\phiXn (t) converges pointwise to\phiX(t) for all t .

\textit{Proof} The real and imaginary parts of eitx are cos(tx) and sin(tx) , which are con-

tinuous and bounded. Therefore E[cos(tXn)]\to E[cos(tX)] and E[sin(tXn)]\to

E[sin(tX)] as n approaches infinity. Therefore E[eitXn ] converges to E[eitX ] as

n\to\infty \hfill $\square$

Using the equivalent conditions in Theorem 14.2, we can show that convergence

in total variation implies weak convergence and thus complete the proof of

Theorem 10.6.

\begin{thmbox}{14.4}
\end{thmbox}

If probability measures \mun,f o r n= 1 ,2,3,... , defined on the measurable

space .(R,\mathcal{B}(R)) converge in total variation distance, then they converge

weakly.

\textit{Proof} Suppose that .(\mun)\infty

n=1 converges in total variation to. \mu, ie.,dTV (\mun,\mu )\to 0

as n\to\infty For any fixed index n, define a measure \nun \coloneqq(\mu n +\mu)/ 2We have

\mun \ll \nun and \mu\ll \nun. By applying Radon-Nikodym theorem (Theorem 13.5), we

can represent \mun by density fn, and \mu by density f, relative to the measure \nun.

Leth: R \to R denote a bounded and continuous function, andhmax be an upper

bound of h(x) over all x. By triangle inequality (Theorem 7.1), we get

\vert\vert\vert

\int

R

h(x) d\mun(x) -

\int

R

h(x) d\mu(x)

\vert\vert\vert =

\vert\vert\vert

\int

R

h(x)fn(x)d\nun(x)

-

\int

R

h(x)f(x)d\nu n(x)

\vert\vert\vert

\leq hmax

\int

R

\vertfn(x) - f( x)\vert d\nun(x).

The proof of Theorem 8.6 can be adapted to probability measures that are absolutely

continuous with respect to any probability measure. Indeed, we can use similar

argument to prove that

dTV (\mun,\mu )= 1

\int

R

\vertfn(x) - f( x)\vert d\nun(x),

which yields

\vert\vert\vert

\int

R

h(x) d\mun(x) -

\int

R

h(x) d\mu(x)

\vert\vert\vert \leq 2hmaxdTV (\mun,\mu ) .

Since dTV (\mun,\mu )converges to 0, we also have .

\int

hd\mu n \to

\int

hd\mu as n \to\infty .

The function h(x) can be any continuous and bounded function in the previous

paragraphs. This proves that probability measure \mun converges weakly to \mu. \hfill $\square$

We note that the converse of Theorem 14.4 does not hold in general. This is

because the total variation distance between a discrete distribution and a continuous

distribution is always equal to 1. For instance, the total variation distance between

a binomial distribution and a normal distribution is equal to 1. However, there

exist examples where binomial distribution converges in distribution to the normal

distribution, despite having a total variation distance of 1 between them. (see

example 10.3.1)
